\begin{proof}[Proof of \cref{obj:lem:staircase-refinement}]
\begin{enumerate}
\item
By \cref{obj:ass:fixed-privacy}, the privacy ratio
\[
r:=e^\varepsilon
\]
satisfies \(r>1\). In the four-record order, the probabilities are
\[
(\pi_{\theta,1},\pi_{\theta,2},\pi_{\theta,3},\pi_{\theta,4})
=
\bigl((1-p)(1-\mu_0),(1-p)\mu_0,
      p(1-\mu_1),p\mu_1\bigr).
\]
Thus \cref{obj:ass:interior-assignment,obj:ass:interior-means} give
\[
\pi_{\theta,j}>0\quad(j=1,\ldots,4),
\qquad
\sum_{j=1}^{4}\pi_{\theta,j}=1.
\]

\item
Define the row measures, their common dominating measure, and the staircase rays by
\[
Q_{\mathrm{row},j}:=Q(\cdot\mid x_j),\qquad
\nu:=\sum_{j=1}^{4}Q_{\mathrm{row},j},\qquad
f_j:=\frac{dQ_{\mathrm{row},j}}{d\nu},
\qquad
b_S(j):=
\begin{cases}
r,&x_j\in S,\\
1,&x_j\notin S.
\end{cases}
\]
Choose nonnegative finite measurable versions of the densities \(f_j\), setting them to zero on any exceptional null set where the Radon--Nikodym derivatives are infinite. The privacy inequalities imply
\[
f_j(z)\le r f_k(z)
\quad\text{for \(\nu\)-almost every \(z\), for every \(j,k\)}.
\]
Indeed, integrating over the measurable set where \(f_j>rf_k\) would contradict
\(Q_{\mathrm{row},j}(A)\le rQ_{\mathrm{row},k}(A)\) if that set had positive \(\nu\)-measure. There are finitely many pairs, so these inequalities hold simultaneously outside one null set.

For every \(z\in Z\), define
\[
m_{\mathrm{base}}(z):=\min_{1\le j\le4}f_j(z),
\qquad
y_j(z):=
\begin{cases}
0,&m_{\mathrm{base}}(z)=0,\\[2mm]
\dfrac{f_j(z)/m_{\mathrm{base}}(z)-1}{r-1},
   &m_{\mathrm{base}}(z)\ne0.
\end{cases}
\]
On the common set where the privacy inequalities hold, an index attaining the minimum gives
\[
0\le m_{\mathrm{base}}(z)\le f_j(z)
   \le r m_{\mathrm{base}}(z).
\]
If \(m_{\mathrm{base}}(z)=0\), all four densities vanish. Otherwise, division by the positive minimum and by \(r-1\) gives \(0\le y_j(z)\le1\). In either case,
\[
f_j(z)=m_{\mathrm{base}}(z)\bigl(1+(r-1)y_j(z)\bigr).
\]

For every subset \(T\subseteq\mathcal X\), including the empty and full subsets, define
\[
\omega_T(z):=
\prod_{j:\,x_j\in T}y_j(z)
\prod_{j:\,x_j\notin T}(1-y_j(z)).
\]
Expanding the products yields
\[
\sum_{T\subseteq\mathcal X}\omega_T(z)=1,
\qquad
\sum_{\substack{T\subseteq\mathcal X\\x_j\in T}}\omega_T(z)=y_j(z),
\]
and consequently
\[
\sum_{T\subseteq\mathcal X}\omega_T(z)b_T(j)
=1+(r-1)y_j(z).
\]
To express this sum using proper nonempty subsets, put
\[
S_\dagger:=\{x_1\},
\qquad
k_{\mathrm{cube}}(z):=
\frac{\omega_\varnothing(z)+r\omega_{\mathcal X}(z)}{r+1},
\]
and, for \(S\in\mathcal S\), define
\[
\lambda_S(z):=
\max\left\{
0,\,
m_{\mathrm{base}}(z)
\left[
\omega_S(z)+k_{\mathrm{cube}}(z)
\bigl(\mathbf1_{\{S=S_\dagger\}}
+\mathbf1_{\{S=\mathcal X\setminus S_\dagger\}}\bigr)
\right]
\right\}.
\]
These functions are measurable and nonnegative on all of \(Z\). On the common set where the privacy inequalities hold, their expressions inside the maximum are nonnegative. Moreover,
\[
b_{S_\dagger}(j)+b_{\mathcal X\setminus S_\dagger}(j)=r+1,
\]
so the added complementary-pair contribution is precisely
\(\omega_\varnothing+r\omega_{\mathcal X}\), the contribution of the empty and full subsets. Therefore
\[
f_j(z)=\sum_{S\in\mathcal S}\lambda_S(z)b_S(j)
\quad\text{for \(\nu\)-almost every \(z\), for every \(j\)}.
\]
The case \(m_{\mathrm{base}}=0\) gives zero on both sides.

\item
Define coefficient measures and weights by
\[
\Lambda_S(A):=\int_A\lambda_S(z)\,d\nu(z),
\qquad
\alpha_S:=\Lambda_S(Z).
\]
Since every ray coordinate is at least one, the preceding decomposition gives
\[
0\le\lambda_S(z)
\le\lambda_S(z)b_S(1)
\le f_1(z)
\quad\text{\(\nu\)-almost everywhere}.
\]
Thus each coefficient is integrable. Integrating the decomposition gives, for every measurable \(A\subseteq Z\),
\[
Q_{\mathrm{row},j}(A)=\sum_{S\in\mathcal S}b_S(j)\Lambda_S(A).
\]
In particular,
\[
\alpha_S\ge0,
\qquad
\sum_{S\in\mathcal S}\alpha_S b_S(j)=Q_{\mathrm{row},j}(Z)=1.
\]
Hence \(\alpha\in\mathcal A_\varepsilon\).

Define the postprocessing measures on the whole pattern set by
\[
K(S,A):=
\begin{cases}
\Lambda_S(A)/\alpha_S,&\alpha_S>0,\\
Q_{\mathrm{row},1}(A),&\alpha_S=0.
\end{cases}
\]
Every row is a probability measure, and the finite pattern domain makes \(K\) a probability kernel. If \(\alpha_S=0\), the nonnegative measure \(\Lambda_S\) has zero total mass and is the zero measure. Consequently, in both cases,
\[
\Lambda_S=\alpha_S K(S,\cdot).
\]
It follows that
\[
Q_{\mathrm{row},j}(A)
=\sum_{S\in\mathcal S}\alpha_S b_S(j)K(S,A)
=\sum_{S\in\mathcal S}Q_\alpha(S\mid x_j)K(S,A),
\]
which proves \(Q=K\circ Q_\alpha\). Indexing these measures and weights by the stated fourteen patterns gives the required channel.

\item
Both information matrices are symmetric: their defining summands and integrands are symmetric outer products. We next identify their quadratic forms.

With \(p\) fixed, define the input score vectors by
\[
\sigma_j:=\frac{\nabla_\theta\pi_{\theta,j}}{\pi_{\theta,j}}
\in\mathbb R^2.
\]
The positive record probabilities make these vectors well defined. Differentiating the pattern masses, with \(p\) and \(\varepsilon\) fixed, gives
\[
h_S(\theta)=\sum_{j=1}^{4}\pi_{\theta,j}b_S(j)>0,
\qquad
g_S
=\sum_{j=1}^{4}b_S(j)\nabla_\theta\pi_{\theta,j}
=\sum_{j=1}^{4}\pi_{\theta,j}b_S(j)\sigma_j.
\]

For an arbitrary direction \(\xi\in\mathbb R^2\), define the release masses, score numerators, and release scores by
\[
\beta_S:=
\sum_{j=1}^{4}\pi_{\theta,j}Q_\alpha(S\mid x_j)
=\alpha_S h_S(\theta),
\]
\[
n_{S,\xi}:=
\sum_{j=1}^{4}\pi_{\theta,j}Q_\alpha(S\mid x_j)
\,\xi^\top\sigma_j
=\alpha_S g_S^\top\xi,
\qquad
\eta_{S,\xi}:=
\begin{cases}
n_{S,\xi}/\beta_S,&\beta_S>0,\\
0,&\beta_S=0.
\end{cases}
\]
If \(\alpha_S=0\), both \(\beta_S\) and \(n_{S,\xi}\) vanish. If \(\alpha_S>0\), cancellation gives
\[
\eta_{S,\xi}=\frac{g_S^\top\xi}{h_S(\theta)},
\qquad
\frac{n_{S,\xi}^{\,2}}{\beta_S}
=\alpha_S\frac{(g_S^\top\xi)^2}{h_S(\theta)}.
\]
Using zero for the contribution of a zero-mass atom, expansion of the quadratic form therefore yields
\[
\sum_{S\in\mathcal S}\beta_S\eta_{S,\xi}^{\,2}
=
\sum_{S\in\mathcal S}
\alpha_S\frac{(g_S^\top\xi)^2}{h_S(\theta)}
=\xi^\top I_\theta(\alpha)\xi.
\]

For the original output, define
\[
f_{\mathrm{out}}(z):=
\sum_{j=1}^{4}\pi_{\theta,j}f_j(z),
\qquad
\mathbf d_\theta(z):=
\sum_{j=1}^{4}\nabla_\theta\pi_{\theta,j}f_j(z),
\]
\[
\zeta_\xi(z):=
\begin{cases}
\dfrac{\xi^\top\mathbf d_\theta(z)}{f_{\mathrm{out}}(z)},
   &f_{\mathrm{out}}(z)>0,\\
0,&f_{\mathrm{out}}(z)=0.
\end{cases}
\]
Here \(f_{\mathrm{out}}\) is the output-law density and
\(\mathbf d_\theta\) is its vector of parameter-derivative densities.
Since the \(\pi_{\theta,j}\) are positive and the \(f_j\) are nonnegative,
\(f_{\mathrm{out}}=0\) implies that every \(f_j=0\), and hence
\(\mathbf d_\theta=0\). Thus, with the same zero-denominator convention,
\[
\xi^\top I_\theta(Q)\xi
=
\int_Z
\frac{(\xi^\top\mathbf d_\theta(z))^2}{f_{\mathrm{out}}(z)}
\,d\nu(z)
=
\int_Z f_{\mathrm{out}}(z)\zeta_\xi(z)^2\,d\nu(z).
\]
The first identity follows by expanding the square and exchanging the finite coordinate sums with the integral.

\item
We derive the exact information-gap identity. The factorization, weighted first by the input probabilities and then by their directional score numerators, gives the following identities of measures:
\[
\sum_{S\in\mathcal S}\beta_S K(S,\cdot)
=f_{\mathrm{out}}\,\nu,
\]
\[
\begin{aligned}
\sum_{S\in\mathcal S}\beta_S\eta_{S,\xi}K(S,\cdot)
&=\sum_{j=1}^{4}
  \pi_{\theta,j}\xi^\top\sigma_j\,Q_{\mathrm{row},j}\\
&=(\xi^\top\mathbf d_\theta)\,\nu.
\end{aligned}
\]
The second identity uses \(\beta_S\eta_{S,\xi}=n_{S,\xi}\), which also holds when \(\beta_S=0\).

The output score is measurable and bounded, because, wherever \(f_{\mathrm{out}}>0\), it is a weighted average of the four directional input scores:
\[
\zeta_\xi(z)
=
\sum_{j=1}^{4}
\frac{\pi_{\theta,j}f_j(z)}{f_{\mathrm{out}}(z)}
\,\xi^\top\sigma_j,
\qquad
|\zeta_\xi(z)|
\le\max_{1\le j\le4}|\xi^\top\sigma_j|.
\]
Its assigned value zero satisfies the same bound when \(f_{\mathrm{out}}=0\). Hence the following integrals are finite. Probability normalization of each row of \(K\), followed by the two measure identities above, gives
\[
\sum_{S\in\mathcal S}\beta_S
\int_Z\eta_{S,\xi}^{\,2}\,K(S,dz)
=\xi^\top I_\theta(\alpha)\xi,
\]
\[
\begin{aligned}
\sum_{S\in\mathcal S}\beta_S\eta_{S,\xi}
\int_Z\zeta_\xi(z)\,K(S,dz)
&=\int_Z(\xi^\top\mathbf d_\theta(z))\zeta_\xi(z)\,d\nu(z)\\
&=\xi^\top I_\theta(Q)\xi,
\end{aligned}
\]
and
\[
\sum_{S\in\mathcal S}\beta_S
\int_Z\zeta_\xi(z)^2\,K(S,dz)
=
\int_Z f_{\mathrm{out}}(z)\zeta_\xi(z)^2\,d\nu(z)
=\xi^\top I_\theta(Q)\xi.
\]
Expanding the square and substituting these three identities proves
\[
\xi^\top\bigl(I_\theta(\alpha)-I_\theta(Q)\bigr)\xi
=
\sum_{S\in\mathcal S}\beta_S
\int_Z\bigl(\eta_{S,\xi}-\zeta_\xi(z)\bigr)^2\,K(S,dz).
\]

\item
Each release mass \(\beta_S\) is nonnegative, and each integral of a square is nonnegative. Therefore, for every \(\xi\in\mathbb R^2\),
\[
\xi^\top\bigl(I_\theta(\alpha)-I_\theta(Q)\bigr)\xi\ge0.
\]
Together with symmetry, this proves
\[
I_\theta(\alpha)-I_\theta(Q)\succeq0.
\]

\item
Fix \(t\in\mathbb R\) and suppose that the two quadratic forms agree in the direction \(v(t)\). The identities above turn this equality into
\[
0=
\sum_{S\in\mathcal S}\beta_S
\int_Z
\bigl(\eta_{S,v(t)}-\zeta_{v(t)}(z)\bigr)^2\,K(S,dz).
\]
Every summand is nonnegative. For any \(S\) with \(\alpha_S\ne0\), nonnegativity of \(\alpha_S\) and positivity of \(h_S(\theta)\) give
\[
\beta_S=\alpha_S h_S(\theta)>0,
\qquad
\eta_{S,v(t)}
=\frac{g_S^\top v(t)}{h_S(\theta)}
=\rho_S(\theta,t).
\]
Its squared residual consequently has zero integral:
\[
\int_Z
\bigl(\rho_S(\theta,t)-\zeta_{v(t)}(z)\bigr)^2\,K(S,dz)=0.
\]
A nonnegative integrable function with zero integral vanishes almost everywhere, so
\[
\zeta_{v(t)}(z)=\rho_S(\theta,t)
\quad\text{for \(K(S,\cdot)\)-almost every \(z\)}.
\]

Now let \(S,T\) be two active patterns. If their projected scores were different, define the measurable set
\[
E_{S,t}:=\{z\in Z:\zeta_{v(t)}(z)\ne\rho_S(\theta,t)\}.
\]
The almost-everywhere score identities for \(S\) and \(T\) would give
\[
K(S,E_{S,t})=0,
\qquad
K(T,Z\setminus E_{S,t})=0.
\]
These equalities establish mutual singularity of the two conditional output measures. Thus, when those measures are not mutually singular, their projected scores must satisfy
\[
\rho_S(\theta,t)=\rho_T(\theta,t),
\]
as claimed.

\end{enumerate}
\end{proof}
